\section{Discussion and Limitations}
\label{sec:discussion}

\paragraph{A hierarchy of evidence.}
The experiments support treating linear availability, generic intervention sensitivity, clinical direction specificity, and mechanism confirmation as different evidentiary levels. A positive probe interval answers whether a restricted readout reliably recovers a label under the patient split and control task. A random/sham pass answers whether the target normal produces an effect beyond matched non-semantic vectors. Clinical-direction comparisons test whether that response belongs preferentially to the named finding; ownership and input closure then test the proposed read--write account. Collapsing these questions turns the original Qwen response into an attribution, whereas the full sequence supplies no registered clinical owner, alias, or closure effect.

\paragraph{Implication for causal interpretation.}
Matched clinical directions are a load-bearing control when a probe normal carries a semantic name. Random vectors test the ambient geometry, and a coordinate permutation tests one construction artifact, but neither represents a plausible competing clinical direction learned with the same readout and data. For the Qwen cell, random/sham controls isolate a non-generic writer but do not identify an Effusion-specific one. The practical rule is therefore to compare semantic directions at the same locus, dose, rows, and endpoint, then test the implied mechanism with multiplicity handled over each fixed family.

\paragraph{Scope of the evidence.}
The registered evidence covers NIH ChestX-ray14, two 7B architectures, the consumed final visual block, Effusion in both models, Edema in LLaVA, six Qwen clinical questions, a paired Consolidation fallback, fixed yes/no prompts, linear readouts, and relative-token probe-normal interventions. It establishes an availability-versus-attribution dissociation for these cells and protocols. The paired decoding intervals span zero, so they leave architecture and concept ordering unresolved. Intermediate loci, nonlinear directions, and other medical image modalities define separate tests rather than extensions of the present estimate.

\paragraph{Data and endpoint limitations.}
ChestX-ray14 labels are mined from reports, with 179 positive Edema test images versus 700 positive Effusion images. The patient split prevents identity leakage but does not remove label noise or acquisition correlations. Mean pooling also does not localize a finding spatially. The behavioural endpoint is the probability assigned to a yes answer under fixed prompts, so the evidence concerns this answer pathway rather than diagnostic accuracy or report quality.

\paragraph{Specificity and closure design.}
The independent supplement tests one discovery-locked Qwen dose against five prospectively fixed clinical normals, and the ownership matrix expands the same comparison to six matching questions. Their maxima are recomputed inside the registered bootstrap families. A broader clinical vocabulary could contain stronger or more closely related directions. In the closure gate, the measured displacement produces $|\alpha|$ of order $10^{-4}$ with no capped pairs; its displacement lower bound crosses zero, and the resulting near-zero gains are not separable from random or clinical controls. The nonconfirmation therefore bounds this registered read--write construction rather than all causal mechanisms in the model.

\paragraph{Reproducibility and ethics.}
Appendix~\ref{app:protocol} records prompts, row-selection hashes, seeds, doses, model revisions, run identifiers, and GPU wall time. Immutable snapshots and terminal receipts bind each result to its environment and inputs. This work evaluates interpretability evidence, not clinical diagnostic performance. Its safety implication is methodological: a decodable feature should be treated as clinically operative only after intervention controls establish direction specificity.
